\documentclass{article}
\usepackage[T1]{fontenc}
\usepackage[utf8]{inputenc}
\begin{document}

This is the opening paragraph of the plain text fixture. It describes the
northbound migration of the quokka population across the veldt, a journey
marked by distinctive foraging patterns and communal nesting behaviour. The
sentences should wrap naturally across the full measure of the text block so
that the PDF inspector's text extraction returns ordinary prose rather than
structured markup. This paragraph deliberately contains no headings, no lists,
and no emphasis commands of any kind, ensuring that it stays plain text from
start to finish. Filler words about weather, soil, and seasonal herbivore
movement fill out the rest of this page so that extraction has real prose to
work with.

\newpage

This second page continues the prose-only fixture. Among the observations
recorded here is the distinctive term \texttt{plaintextsentinel-foxtrot},
which is used only once in this document so that automated tests can assert
that it appears verbatim in the extracted content. The remainder of this
paragraph is deliberately generic filler prose about rainfall, drainage, and
the gradual retreat of the frost line. No structure is introduced here either;
the text remains a sequence of unformatted paragraphs with nothing but a page
break separating it from its neighbours.

\newpage

The third and final page closes out the plain text document. It carries the
unique marker \texttt{plaintextsentinel-golf} so that a test can verify that
text spanning the third page is captured. A closing remark reiterates that
this fixture is intentionally plain: paragraphs only, with pagination
controlled by explicit page breaks and no visual formatting beyond the
article class defaults. With this closing sentence the plain text fixture
reaches its end.

\end{document}
